\section{Sesgo: más sutil de lo que imaginas}

\subsection{Sesgo de género en consultas SQL}

El expositor compartió un caso vivido en persona para ilustrar la sutileza del sesgo en la IA. Mientras trabajaba en un proyecto de genealogía, le pidió a un chatbot que generara una consulta SQL para buscar a todos los ancestros. El código devuelto tenía un error crucial: en la consulta recursiva, una línea \textbf{solo verificaba si el father era ancestro y omitía a la mother}. Aunque el resto del código sí trataba el lado materno, este paso recursivo específico lo dejó fuera.

\begin{warningbox}{¿Bug o sesgo? Una pregunta difícil de responder}
El expositor discutió el problema con su esposa: ¿esto era en realidad un \textbf{error de programación} (Bug) o un \textbf{sesgo de género} (Gender Bias)? Si en los datos de entrenamiento los ejemplos de código de genealogía son mayoritariamente patrilineales, el modelo terminaría «aprendiendo» ese patrón de preferencia. El problema central es que \textbf{no puedes preguntarle a un LLM si tiene sesgo}: no cuenta con una capacidad confiable de autoexamen. Más adelante en la conversación, el chatbot sí corrigió el error, pero la salida inicial ya había revelado el sesgo latente.
\end{warningbox}

\subsection{El sesgo no se resuelve prohibiendo palabras}

El expositor señaló un malentendido común: la gente cree que se puede eliminar el sesgo de un LLM bloqueando ciertas palabras, por ejemplo, impidiendo que aparezcan en la salida términos relacionados con la raza. Pero el sesgo es mucho más profundo que palabras específicas: puede estar implícito en la estructura de las oraciones, en la selección de temas, en qué detalles se incluyen o se omiten, y en muchos otros niveles.

Un caso clásico: los dispensadores automáticos de jabón, cuyos sensores estaban calibrados solo para piel clara, no liberaban jabón para usuarios de piel oscura. Esto no involucra ninguna «palabra sesgada», sino un sesgo sistémico en el diseño técnico mismo. En los modelos de lenguaje de gran escala existen sesgos análogos, presentes de forma más encubierta tanto en el código como en el texto que generan.

\subsection{Resumen de la sección}

El sesgo en la IA es más sutil y está más arraigado de lo que aparenta a simple vista. Puede manifestarse como preferencia patrilineal en el código, como sesgo demográfico en los resultados de recomendación, o como estereotipos en la generación de lenguaje. Una estrategia simple de prohibir palabras no resuelve este problema; se necesitan pruebas y evaluaciones sistemáticas para detectar y mitigar el sesgo.

